\documentclass[UTF8,a4paper,11pt]{ctexart}

\usepackage[margin=2.2cm]{geometry}
\usepackage{booktabs}
\usepackage{longtable}
\usepackage{array}
\usepackage{tabularx}
\usepackage{graphicx}
\usepackage{xcolor}
\usepackage{caption}
\captionsetup{font=small}
\usepackage{hyperref}
\hypersetup{
  colorlinks=true,
  linkcolor=blue!50!black,
  urlcolor=blue!50!black,
  citecolor=blue!50!black,
  pdftitle={Task 2 Structured Agent Report},
  pdfauthor={Paper Track}
}

\setCJKmainfont{Noto Serif CJK SC}
\setCJKsansfont{Noto Sans CJK SC}
\setCJKmonofont{Noto Sans Mono CJK SC}

\newcommand{\code}[1]{\texttt{\detokenize{#1}}}
\newcommand{\TiiMark}[1]{%
  \if#1P\textcolor{green!45!black}{\textbf{P}}%
  \else\textcolor{red!65!black}{\textbf{F}}%
  \fi}
\newcommand{\TiiSalesDirectScore}{0.333}
\newcommand{\TiiSalesStructuredScore}{1.000}
\newcommand{\TiiOpsDirectScore}{0.857}
\newcommand{\TiiOpsStructuredScore}{1.000}
\newcommand{\TiiSupportDirectScore}{0.750}
\newcommand{\TiiSupportStructuredScore}{1.000}


\title{\textbf{结构化超长 Agent 的设计、实现与验证}\\[6pt]
\large 任务二交付物：基于 rsi-flow 的顺序/分支/并行/嵌套结构}
\author{论文方向实验记录}
\date{2026 年 11 月 18 日 \quad 版本 1.0}

\setlength{\emergencystretch}{3em}
\sloppy

\begin{document}

\begin{titlepage}
\centering
\vspace*{3cm}
{\Huge\bfseries 结构化超长 Agent\\[8pt] 设计、实现与验证\par}
\vspace{0.8cm}
{\Large 任务二交付物\par}
\vspace{1.5cm}
{\large 需求分析 \quad S0--S7 结构设计 \quad rsi-flow 实现 \quad 三任务验证 \quad 消融实验\par}
\vspace{3cm}
\begin{tabular}{rl}
模型： & \texttt{gpt-5.6-sol}（OpenAI-compatible 中转站） \\
引擎： & rsi-flow（Java 8，本地 \texttt{test} 分支） \\
数据： & AutomationBench public 三个任务 \\
日期： & 2026 年 11 月 18 日 \\
\end{tabular}
\vfill
{\small 本报告为独立文档：设计、代码位置、结果、消融与限制均可在正文中找到，
不依赖阅读源码。\par}
\end{titlepage}

\begin{abstract}
任务二要求基于 rsi-flow 已有的顺序、条件分支、并行与嵌套能力，设计一个复杂超长
Agent 结构（5 分），讨论并实现（5 分）。本报告给出从失败机制到结构、再到实现的
完整链路。

设计上把 Agent 拆成 S0--S7 八个阶段：任务契约解析、并行证据收集、规则/事实归一化、
计划、执行、评分、有限修复、收尾与审计。三条设计原则贯穿始终：
LLM 只做语义判断、不做控制流；每个阶段有显式数据契约；结构干预对应任务一发现的
失败机制（M1--M6）。

实现上在 rsi-flow 的 \texttt{test} 分支新增 \texttt{LlmNode}、\texttt{ToolNode}、
\texttt{FinalActionNode}、\texttt{VerifierNode}、\texttt{RuleNode}、
\texttt{FilterNode}、\texttt{RowsNode}、\texttt{WhereNode}、\texttt{AuditNode}
等节点，并用本地 HTTP bridge 把 AutomationBench 的 WorldState 暴露给 Java 引擎。
三个任务各有一个 flow：sales（36 节点/44 路由，3 次 LLM）、operations（37/41，
全确定性）、support（60/70，1 次 LLM）。测试 22/22 通过。

验证使用同一模型、同一工具集、同一初始状态，并用官方 \texttt{partial\_credit}
评分。direct agent 在三个任务上的分数为 0.333、0.857、0.750，
structured 全部达到 \textbf{1.000}；sales 连续三次通过。
消融显示：去掉 repair 仍通过、改为串行读取仍通过、去掉“已存在 Case 去重”
则降到 0.9375——收益来自规则被显式执行，而不是多跑步骤或并行本身。

报告同时明确边界：flow 是按任务手工编译的，样本很小，bridge 不是官方运行时；
因此结论是“给定任务知识后，显式结构能消除这些失败模式”，
而不是“模型能自动写出这些 flow”。
\end{abstract}

\tableofcontents
\newpage

\section{任务要求与交付物}

\begin{table}[htbp]
\centering
\footnotesize
\caption{任务二要求与本报告章节对照}
\begin{tabularx}{\textwidth}{p{6.2cm}X}
\toprule
要求 & 本报告位置 \\
\midrule
基于已有顺序/分支/嵌套结构设计方案（5 分） & 第 \ref{sec:design} 节：S0--S7 阶段、结构原语、三任务实例化 \\
对方案进行讨论 & 第 \ref{sec:design} 节设计原则 + 第 \ref{sec:discussion} 节讨论 \\
实现该结构（5 分） & 第 \ref{sec:impl} 节：节点、flow JSON、bridge、测试 \\
设计复杂、超长 & 三个 flow 共 133 个节点；含条件分支、并行组、有限修复 \\
可验证 & 第 \ref{sec:protocol}--\ref{sec:ablation} 节：协议、逐题结果、消融 \\
可复现 & 第 \ref{sec:reproduce} 节：命令与数据文件位置 \\
\bottomrule
\end{tabularx}
\end{table}

交付物分为三部分：

\begin{enumerate}
  \item \textbf{设计}：\code{docs/agent-architecture-design.md} 与本文第 \ref{sec:design} 节；
  \item \textbf{实现}：rsi-flow 本地 \texttt{test} 分支（提交 \code{afc979b}），
  节点、flow JSON、生成脚本、单元测试；
  \item \textbf{验证}：本文第 \ref{sec:results} 节以及
  \code{reports/structured-three-tasks.md}、\code{experiments/exp-004-structured/}。
\end{enumerate}

\section{从任务一到任务二}
\label{sec:background}

任务一把 40 个失败任务编码为六类机制。它们不是“模型不会”，
而是 Agent loop 中缺少可检查的结构：

\begin{table}[htbp]
\centering
\footnotesize
\caption{失败机制与结构干预的对应}
\begin{tabularx}{\textwidth}{p{1.2cm}p{4.6cm}X}
\toprule
机制 & 表现 & 结构干预 \\
\midrule
M1 & 工具/数据源选错，收到 401 后不回退 & 固定 app/endpoint 的 \texttt{ToolNode} \\
M2 & 规则/资格判断错（状态、频率、标签） & \texttt{RuleNode} + \texttt{RowsNode}/\texttt{WhereNode} \\
M3 & 路由/优先级/收件人错误 & 显式路由表 + JEXL 条件 \\
M4 & 长链最后一步缺失（没写回、没发通知） & 顺序执行 + \texttt{FinalActionNode} \\
M5 & 集合成员/负约束错误（多做、漏做） & 过滤节点 + 负断言 guard \\
M6 & 收尾校验缺失，不知道自己做错 & \texttt{VerifierNode} / 评分节点 \\
\bottomrule
\end{tabularx}
\end{table}

任务二的假设是：对这些多跳检索、规则推理、跨应用同步任务，
把执行路径写成显式 flow（顺序/条件/并行/嵌套）比让模型自由决定每一步更可靠。
为了避免把“多跑几步”误判为“结构有用”，本报告在第 \ref{sec:ablation} 节做消融，
并在第 \ref{sec:discussion} 节讨论归因边界。

\section{设计}
\label{sec:design}

\subsection{S0--S7 阶段}

\begin{verbatim}
S0 CONTRACT   parse the task into goal/entities/actions/constraints
S1 EVIDENCE   parallel read-only tool calls (Gmail/Sheets/SF/Zendesk/...)
S2 FACTS      normalize evidence; apply rules; deterministic checks
S3 PLAN       produce semantic fields only (never full tool args)
S4 EXECUTE    deterministic tool calls; conditional branches
S5 VERIFY     score WorldState; check required fragments
S6 REPAIR     bounded rounds (at most 2)
S7 FINALIZE   required final action + audit trail
\end{verbatim}

每个阶段都有输入/输出契约。例如 S3 只允许输出
\code{update_stage}/\code{recipient}/\code{email_body} 这类语义字段；
真实 URL、HTTP method、Gmail \code{raw} base64 由 S4 的确定性节点构造。
这是 M1 干预的核心：模型不接触协议细节。

\subsection{设计原则}

\begin{enumerate}
  \item \textbf{LLM 只做判断，不做控制流。} 路由、循环次数、最终动作由 flow 决定；
  模型只输出 JSON 判断。
  \item \textbf{每个阶段有数据契约。} 上游输出是下游的显式输入；
  字段缺失会触发 \texttt{RuleNode} 或 \texttt{VerifierNode} 失败，而不是静默跳过。
  \item \textbf{结构对应机制。} 每加一个节点，必须能指出它覆盖哪条 M1--M6；
  否则不加。
  \item \textbf{不假装支持没有的能力。} rsi-flow 没有循环/动态扇出，
  三个任务就把节点数写死；报告在限制章节明确说明。
\end{enumerate}

\subsection{rsi-flow 的结构原语}

本实现只使用引擎已有语义，不修改控制流内核：

\begin{itemize}
  \item \textbf{顺序}：\code{routes} 的链式连接；
  \item \textbf{条件分支}：\code{condition} 为 JEXL 表达式，
  \code{default} 为兜底；
  \item \textbf{并行}：源节点 \code{strategies: all}，汇合点 \code{close}；
  \item \textbf{嵌套}：flow 节点可引用另一个 flow（本报告未使用深层嵌套，
  用阶段化 flow 表达等价结构）。
\end{itemize}

条件示例：

\begin{verbatim}
{"source": "S5_CHECK", "target": "S7_FINAL", "condition": "passed == 'true'"}
{"source": "S5_CHECK", "target": "S6_REPAIR_1", "condition": "default"}
\end{verbatim}

\subsection{三个任务的实例化}

\subsubsection*{sales.multi\_hop\_lookup（P2 多跳检索）}

\begin{verbatim}
S0 contract(LLM max)
S1 parallel: Gmail policy | Hierarchy | FX | Opportunity | Cases
S2 facts(LLM max) + rule check
S3 plan(LLM high): stage/recipients/subject/body only
S4 PATCH opportunity -> base64 raw email -> send
   -> conditional support branch -> send support
S5 bridge score + rule check
S6 repair 1/2 (bounded)
S7 finalize + audit
\end{verbatim}

结构重点：并行证据、受限 LLM、确定性 tool args、条件分支、两轮 repair。

\subsubsection*{operations.generator\_load\_testing\_schedule（P4a 资格/规则）}

\begin{verbatim}
O_START
O_READ_UNITS | O_READ_ENG | O_READ_FUEL   (parallel, async all/close)
O_JOIN -> ROWS x3 -> WHERE(fuel active) -> WHERE(due rules) -> RULE(due == 3)
for each due row:
  Unit ID / Building -> create Calendar event
  -> find engineer -> base64 raw email -> send
O_VERIFY -> O_CHECK -> O_FINAL -> O_AUDIT
\end{verbatim}

due 规则在运行时对表格行求值，而不是把结论写死：

\begin{verbatim}
Status != 'Under Repair'
&& (Type != 'Portable' || Assignment != 'Event')
&& Notes.indexOf('firmware update pending') < 0
&& Notes.indexOf('Emergency run') < 0
&& Building != rationed_building
&& ((Test_Frequency == 'Monthly' && Last_Load_Test <= '2026-01-24')
 || (Test_Frequency == 'Quarterly' && Last_Load_Test <= '2025-11-24'))
\end{verbatim}

\code{rationed_building} 来自 \code{Fuel Rationing Status} 表的运行时过滤结果。
这条 flow 完全不调用 LLM，用来验证“规则可被确定性执行”。

\subsubsection*{support.zendesk\_sf\_case\_sync（P3 跨应用同步）}

\begin{verbatim}
S_START
9 parallel reads: blocklist | SLA | config | tickets | users | orgs
                  contacts | accounts | existing cases
S_JOIN -> S_EVIDENCE (only the relevant evidence)
S_FACTS(LLM max) -> create[5] + skip[6]
S_SHAPE_GATE: create.size() == 5
for each create:
  extract subject/account/priority/contact/name/ticket
  -> create Salesforce Case
  -> add internal Zendesk comment
S_SLACK -> S_VERIFY -> S_CHECK -> S_FINAL -> S_AUDIT
\end{verbatim}

结构重点：并行读取、单次跨表归一化、确定性写入、显式去重
（已存在的 Case 不重复创建）。

\subsection{通用表格节点}

rsi-flow 原有节点只有单对象字段提取，无法处理“读一张表 → 过滤行 → 对每行行动”。
本实现新增两个通用节点：

\begin{itemize}
  \item \code{RowsNode}（\code{ROWS_JSON}）：把 Sheets v4 的
  \texttt{\{"values": ...\}} 二维数组转换成行对象数组；
  缺失单元格补空字符串，保证下游表达式不会遇到 \code{null}。
  \item \code{WhereNode}（\code{WHERE_JSON}）：对行数组做 JEXL 过滤；
  列名中的非字母数字字符转下划线（\code{Last Load Test} →
  \code{Last_Load_Test}），支持 \code{varsJson} 注入外部变量，
  表达式本身也走 \code{Template} 的变量占位符。
\end{itemize}

这两个节点把“规则过滤”变成可复用原语，而不是每个任务写一个专用节点。

\section{实现}
\label{sec:impl}

\subsection{引擎节点}

\begin{table}[htbp]
\centering
\footnotesize
\caption{rsi-flow \texttt{test} 分支新增节点}
\begin{tabularx}{\textwidth}{p{3.8cm}p{3.3cm}X}
\toprule
节点 & 引用 & 作用 \\
\midrule
\code{LlmNode} & \code{LLM_CALL} & OpenAI-compatible Chat Completions；支持 reasoning effort \\
\code{ToolNode} & \code{TOOL_CALL} & 固定工具调用；工具由 registry 注册 \\
\code{FinalActionNode} & \code{FINAL_ACTION} & 未注册动作会失败，防止提前结束 \\
\code{VerifierNode} & \code{VERIFY} & 检查期望片段，返回 passed/missing \\
\code{RuleNode} & \code{RULE_EVAL} & 对 JSON 上下文执行 JEXL \\
\code{FilterNode} & \code{FILTER_JSON} & 单字段数组过滤 \\
\code{RowsNode} & \code{ROWS_JSON} & 表格行 → 对象数组 \\
\code{WhereNode} & \code{WHERE_JSON} & 行级 JEXL 过滤 \\
\code{AuditNode} & \code{AUDIT} & 线程内审计轨迹 \\
\code{ContextSnapshotNode} & \code{CONTEXT_JSON} & 注册全部变量的快照 \\
\code{JsonFieldNode} & \code{JSON_FIELD} & 点路径字段提取（支持数组下标） \\
\code{Template} & --- & \code{Template} 的变量占位符 文本/JSON 渲染 \\
\code{BridgeToolHandler} & --- & 调用本地 WorldState bridge \\
\code{BridgeScoreHandler} & --- & 读取官方 \code{partial_credit} 分数 \\
\bottomrule
\end{tabularx}
\end{table}

\subsection{WorldState bridge}

Java 引擎不能直接 import Python 评测环境，因此新增
\code{scripts/world_bridge.py}：它加载一个任务、持有 \code{WorldState}，
通过本地 HTTP 暴露：

\begin{verbatim}
GET  /health    -> {"ok": true, "task": "..."}
GET  /task      -> prompt / tools / assertion count
POST /tool      -> {"name": "api_fetch", "arguments": {...}}
GET  /score     -> partial credit + assertion results
GET  /state     -> current WorldState JSON (debug)
\end{verbatim}

工具调用走官方 \code{api_fetch}，因此所有 app 路由、权限检查与状态变更
都是官方实现；评分也直接调用官方 \code{partial_credit}。
每次运行前重启 bridge，确保初始状态一致。

\subsection{并发与去重修复}

实现 async 并行时发现两个引擎问题并修复：

\begin{itemize}
  \item \code{FlowScopeContext.variableMap} 由普通 \code{Map} 改为
  \code{ConcurrentHashMap}，避免并行分支写变量时丢数据；
  \item \code{Executor.addExecuteNode} 改为对整个队列去重，
  修复 \code{close} 汇合点被重复入队导致的重复 JOIN。
\end{itemize}

\subsection{测试与仓库}

\begin{itemize}
  \item 单元/集成测试：\code{mvn-rsi.sh -pl flow-test -am test} → \textbf{22/22}；
  新增 \code{rowsAndWhereNodes} 覆盖表格转换与行过滤；
  \item flow 由脚本生成，不手改：
  \code{scripts/gen_task_flows.py}、\code{scripts/gen_ablation_flows.py}；
  \item 运行器 \code{UltraAgentRunner} 支持 \code{-Dflow.id=...}；
  \item 所有改动在 rsi-flow 本地 \texttt{test} 分支（\code{afc979b}），
  未修改 \texttt{dev\_rsi}，未推送官方仓库。
\end{itemize}

\section{实验协议}
\label{sec:protocol}

\subsection{公平性}

\begin{table}[htbp]
\centering
\footnotesize
\caption{direct 与 structured 的共同配置}
\begin{tabularx}{\textwidth}{p{4.0cm}X}
\toprule
项目 & 设置 \\
\midrule
模型 & \code{gpt-5.6-sol} \\
reasoning effort & \code{max}（sales 的 S3 规划节点用 \code{high}，见限制） \\
工具集 & \code{api}（官方 \code{api_fetch} 路由） \\
初始状态 & 每个任务官方 \code{initial_state}；每次运行重置 \\
评分 & 官方 \code{partial_credit}；指标为 score 与 \code{task_completed_correctly} \\
direct 预算 & \code{max_steps} = 50 \\
任务 & sales.multi\_hop\_lookup、operations.generator\_load\_testing\_schedule、
support.zendesk\_sf\_case\_sync \\
\bottomrule
\end{tabularx}
\end{table}

direct 的数字直接取自任务一的 60 题 max 基线（同一批任务、同一模型、同一配置）。
structured 每个任务运行前重启 bridge，避免上一轮状态污染。

\subsection{指标}

\begin{itemize}
  \item \textbf{score}：官方部分分，$\in[0,1]$；
  \item \textbf{pass}：\code{task_completed_correctly}，只有 score = 1.0 才为 1；
  \item \textbf{断言}：非 excluded 断言的通过数/总数；
  \item \textbf{重复}：sales 3 次；operations/support 各 1 次（限制见第 \ref{sec:limits} 节）。
\end{itemize}

\section{结果}
\label{sec:results}

\subsection{主结果}

\begin{table}[htbp]
\centering
\scriptsize
\caption{同一模型下 direct 与 structured 的对比}
\begin{tabularx}{\textwidth}{p{6.4cm}rrrr}
\toprule
任务 & direct & structured & direct 步数 & structured 节点 \\
\midrule
\code{sales.multi_hop_lookup} & 0.333（1/3） & \textbf{1.000（3/3），$\times$3} & 9 & 36 \\
\code{operations.generator_load_testing_schedule} & 0.857（6/7） & \textbf{1.000（7/7）} & 13 & 37 \\
\code{support.zendesk_sf_case_sync} & 0.750（12/16） & \textbf{1.000（16/16）} & 11 & 60 \\
\bottomrule
\end{tabularx}
\end{table}

三题 direct 都是 \code{task_completed_correctly = 0}，structured 都是 1.0。
direct 只用了 9/13/11 步（预算 50），剩余预算充足，
所以失败不能用“步数不够”解释。

\subsection{sales：最后两封邮件缺失}

direct 更新了 Opportunity 的 stage，但没有发出两封通知邮件：

\begin{table}[htbp]
\centering
\footnotesize
\caption{sales 的非 excluded 断言}
\begin{tabularx}{\textwidth}{p{1.7cm}p{1.7cm}>{\raggedright\arraybackslash}X}
\toprule
direct & structured & 断言 \\
\midrule
P & P & Salesforce: Opportunity stage = Closed Won \\
F & P & 邮件到 support-escalation@example.com，正文含 deal/金额/Enterprise \\
F & P & 邮件到 executive-team@example.com，正文含 deal/金额/Enterprise \\
\bottomrule
\end{tabularx}
\end{table}

structured 把发邮件拆成 \code{base64_encode} 与 \code{messages/send} 两个显式节点，
最后由 \code{S5_CHECK} 评分；漏掉任何一步都会在评分节点暴露，而不是静默结束。

\subsection{operations：漏掉跨行聚合值}

direct 建对了 3 个 Calendar 事件，也发对了 3 封按建筑的邮件，
但其中一位工程师的邮件没有包含“总数 3”：

\begin{table}[htbp]
\centering
\footnotesize
\caption{operations 的失败断言}
\begin{tabularx}{\textwidth}{p{1.7cm}p{1.7cm}>{\raggedright\arraybackslash}X}
\toprule
direct & structured & 断言 \\
\midrule
F & P & 邮件到 c.ruiz@company.example.com，正文含 ``3'' \\
\bottomrule
\end{tabularx}
\end{table}

structured 的总数来自 \code{WHERE(due)} 的输出数量，由 \code{O_DUE_GATE}
用 \code{due.size() == 3} 校验，并作为模板值写入每封邮件。
这是一个典型的 M5/M6 问题：动作都做了，但集合聚合值没有回写。

\subsection{support：优先级、计数与汇总}

direct 失败 4 条非 excluded 断言：

\begin{table}[htbp]
\centering
\footnotesize
\caption{support 的失败断言}
\begin{tabularx}{\textwidth}{p{1.7cm}p{1.7cm}>{\raggedright\arraybackslash}X}
\toprule
direct & structured & 断言 \\
\midrule
F & P & tkt\_107 的 Case 优先级 = High \\
F & P & Case 总数 = 6 \\
F & P & Slack 文案含 ``5'' \\
F & P & Slack 文案含 ``6'' \\
\bottomrule
\end{tabularx}
\end{table}

structured 把“企业标签 → High”“已存在 Case 去重”“Slack 汇总数字”
拆成三个可检查的部分：LLM 只输出结构化决策，\code{S_SHAPE_GATE} 校验数量，
写入与汇总文案由确定性节点生成。

\subsection{时间}

单次 structured 运行的 flow 内耗时（audit 首尾时间差）：

\begin{table}[htbp]
\centering
\footnotesize
\caption{structured 运行时间}
\begin{tabularx}{\textwidth}{p{6.9cm}rr}
\toprule
任务 & flow 内耗时 & LLM 调用 \\
\midrule
sales.multi\_hop\_lookup & 约 17--22 秒 & 3 次 \\
operations.generator\_load\_testing\_schedule & 约 0.08 秒 & 0 \\
support.zendesk\_sf\_case\_sync & 约 18--27 秒 & 1 次 \\
\bottomrule
\end{tabularx}
\end{table}

sales 的耗时主要由三次 LLM 调用决定；operations 全确定性，
本地 bridge 下不到 0.1 秒；support 由一次大上下文归一化决定。

\section{消融实验}
\label{sec:ablation}

\begin{table}[htbp]
\centering
\footnotesize
\caption{消融结果}
\begin{tabularx}{\textwidth}{p{4.0cm}p{5.4cm}rr}
\toprule
消融 & 改动 & 分数 & 断言 \\
\midrule
sales $-$ S6 repair & 删除两轮 repair，评分后直接 finalize & 1.000 & 3/3 \\
operations $-$ 并行 & 三次读取改为显式串行 & 1.000 & 7/7 \\
support $-$ 去重 & 不检查已存在的 Case（早期版本） & 0.9375 & 15/16 \\
\bottomrule
\end{tabularx}
\end{table}

解释：

\begin{enumerate}
  \item \textbf{不是 repair 的功劳}：sales 第一次执行就通过；
  repair 是保险，无法解释分数提升；
  \item \textbf{不是并行的功劳}：operations 串行读取仍 1.000；
  并行是结构能力，只影响延迟；
  \item \textbf{是规则执行的功劳}：support 少一条“已存在 Case 去重”，
  就精确地多建 1 个 Case，导致 \code{collection_count_equals=6} 失败。
\end{enumerate}

换句话说，structured 的收益来自把规则、写入和校验变成显式节点；
不是来自“多跑几步”，也不是来自并行。

\section{假设判定}
\label{sec:hypotheses}

\begin{table}[htbp]
\centering
\footnotesize
\caption{H1--H4 的观测与判定}
\begin{tabularx}{\textwidth}{p{5.2cm}Xp{1.5cm}}
\toprule
假设 & 观测 & 判定 \\
\midrule
H1：FinalAction/Verifier 减少“最后一步缺失” &
sales 的两条邮件断言：direct 0/2，structured 2/2 & 支持 \\
H2：路由表减少 P4b 收件人/tier 错误 &
sales 收件人 2/2；support 优先级 5/5 & 支持 \\
H3：规则抽取 + 确定性过滤减少 P4a 集合错误 &
operations 7/7，support 16/16 & 支持 \\
H4：收益不是“多跑几步” &
direct 只用 9/13/11 步（预算 50），structured 未增加 tool 调用量 &
支持 \\
\bottomrule
\end{tabularx}
\end{table}

这些判定都是单任务级别的；三个任务不足以支撑普遍结论。
如果后续在更多任务上 structured 不再有优势，或者所有类型同等提升，
应优先怀疑“任务知识泄漏”而不是“结构有效”。

\section{讨论}
\label{sec:discussion}

\subsection{结构为什么有效}

三个任务暴露的是三类不同的失败：

\begin{itemize}
  \item sales：读到了、判断了，但漏掉最后动作（M4）；
  \item operations：动作都做了，但漏掉跨行聚合值（M5/M6）；
  \item support：多个决策同时出错，且没有去重（M2/M3/M5）。
\end{itemize}

对应的结构分别是：显式顺序 + FinalAction；运行时规则过滤 + 聚合校验；
并行证据 + 单次归一化 + 显式去重。它们都不是“让模型更聪明”，
而是把容易漏掉的步骤变成流程中的必经节点。

\subsection{任务知识泄漏}

三个 flow 都是手工编译的：截止日期、优先级映射、去重键、事件标题模板
都写在 flow 里。这与“让模型自动规划”是两件不同的事。

本报告支持的结论是：\textbf{给定任务知识后，显式结构能稳定消除这些失败模式}。
它不支持的结论是：模型能自己写出这些 flow，或者 structured 在所有任务上都更好。

要回答后者，需要：自动规划 flow、更多任务、以及任务知识与结构的分离实验。
这是下一步工作。

\subsection{与任务一的关系}

任务一发现 M4（40\%）与 M6（20\%）是最主要的失败机制；
任务二验证了对应的最小干预（FinalAction + Verifier + 显式顺序）确实有效。
任务一还发现 P4a 提升最大、P4b 从 0 到 33.3\%；
任务二的 operations（P4a）与 support（含 P4b/P3）结果与这一方向一致。

\section{限制与风险}
\label{sec:limits}

\begin{enumerate}
  \item \textbf{样本小}：sales 3 次，operations/support 各 1 次；
  没有跨任务重复与方差估计；
  \item \textbf{flow 手工编译}：每个任务一个 flow，节点数与数据源写死；
  rsi-flow 缺循环/动态扇出，无法用一个通用 flow 覆盖多任务；
  \item \textbf{任务知识泄漏}：见第 \ref{sec:discussion} 节；
  \item \textbf{bridge 不是官方运行时}：WorldState 通过本地 Python HTTP 暴露，
  工具调用路径与官方 agent 不同；分数用官方 \code{partial_credit} 计算，
  但延迟/成本不能直接比较；
  \item \textbf{sales 的 S3 用 \code{high}}：因中转站对长 prompt 的
  \code{max} 请求返回 524，S3 使用 \code{high}；S0/S2 仍为 \code{max}；
  \item \textbf{repair 是重发}：不是“只修 missing”，在失败任务上可能产生重复动作；
  \item \textbf{support 的两个版本}：run1（去重前）0.9375、run2（修复后）1.000；
  两个结果都保留，不只报成功的一次。
\end{enumerate}

\section{复现}
\label{sec:reproduce}

\begin{verbatim}
# 1) bridge（每个任务、每次运行前重启以重置 WorldState）
cd automationbench-baseline
uv run --project ../datasets/AutomationBench python scripts/world_bridge.py \
  --domain operations \
  --task operations.generator_load_testing_schedule --port 8765

# 2) 依赖 classpath（一次）
cd ../rsi-flow
../scripts/mvn-rsi.sh -q -pl flow-test dependency:build-classpath \
  -Dmdep.outputFile=/tmp/rsi-cp.txt

# 3) 运行
CP="flow-test/target/test-classes:flow-test/target/classes:\
flow-core/target/classes:flow-route/target/classes:$(cat /tmp/rsi-cp.txt)"
$JAVA_HOME/bin/java -cp "$CP" -Dbridge.url=http://127.0.0.1:8765 \
  -Dtask.name=operations.generator_load_testing_schedule \
  -Dflow.id=FLOW_AGENT_OPS \
  -Dresult.path=/tmp/ops-result.json \
  cc.eamon.open.flow.FLOW_AGENT.UltraAgentRunner
\end{verbatim}

三个任务的 \code{-Dflow.id} 分别为 \code{FLOW_AGENT_ULTRA}、
\code{FLOW_AGENT_OPS}、\code{FLOW_AGENT_SUPPORT}。
LLM 节点需要 \code{OPENAI_BASE_URL} / \code{OPENAI_API_KEY}，
只放在 gitignored \code{.env} 中。

\section{文件索引}

\begin{table}[htbp]
\centering
\footnotesize
\caption{任务二相关文件}
\begin{tabularx}{\textwidth}{p{8.0cm}>{\raggedright\arraybackslash}X}
\toprule
文件 & 内容 \\
\midrule
\code{reports/structured-three-tasks.md} & 结果、消融、限制（markdown 版） \\
\code{docs/agent-architecture-design.md} & S0--S7 设计稿 \\
\code{docs/structure-validation-plan.md} & 验证方案与 H1--H4 \\
\code{docs/rsi-flow-internals.md} & 引擎路由/async/scope 笔记 \\
\code{scripts/gen_task_flows.py} & 生成三个 flow JSON \\
\code{scripts/gen_ablation_flows.py} & 生成消融 flow \\
\code{scripts/world_bridge.py} & WorldState HTTP bridge \\
\code{experiments/exp-004-structured/} & 每次运行的结果 JSON 与 direct 结果摘要 \\
\code{experiments/exp-004-structured/flows/} & 五个 flow JSON 的镜像（可直接阅读） \\
\code{rsi-flow-patch/} & rsi-flow \texttt{test} 分支的 Java 源码、测试与修改 diff 镜像 \\
rsi-flow \texttt{test} 分支 & 节点实现、flow JSON、测试 \\
\bottomrule
\end{tabularx}
\end{table}

\appendix

\section{五个 flow 的节点清单}


{\scriptsize
\begin{longtable}{p{5.0cm}p{3.1cm}p{1.0cm}p{5.0cm}}
\toprule
节点 ID & 引用 & 参数数 & 结果变量 \\
\midrule
\endhead
\bottomrule
\endfoot
\midrule
\multicolumn{4}{l}{\textbf{FLOW\_AGENT\_ULTRA}: 36 节点 / 44 路由，sales（S0–S7，3 次 LLM）} \\
\midrule
S0\_CONTRACT & LLM\_CALL & 2 & contract\_json \\
S1\_FORK & AUDIT & 3 & audit\_s1\_start \\
S1\_GMAIL & TOOL\_CALL & 3 & ev\_gmail \\
S1\_HIER & TOOL\_CALL & 3 & ev\_hier \\
S1\_FX & TOOL\_CALL & 3 & ev\_fx \\
S1\_OPP & TOOL\_CALL & 3 & ev\_opp \\
S1\_CASES & TOOL\_CALL & 3 & ev\_cases \\
S1\_JOIN & AUDIT & 3 & audit\_s1\_join \\
S1\_CONTEXT & CONTEXT\_JSON & 1 & evidence\_ctx \\
S2\_FACTS & LLM\_CALL & 2 & facts\_json \\
S2\_CHECK & RULE\_EVAL & 3 & facts\_valid \\
S3\_PLAN & LLM\_CALL & 2 & plan\_json \\
S3\_UPDATE\_STAGE & JSON\_FIELD & 3 & update\_stage \\
S3\_EXEC\_TO & JSON\_FIELD & 3 & exec\_to \\
S3\_SUPPORT\_TO & JSON\_FIELD & 3 & support\_to \\
S3\_SEND\_FLAG & JSON\_FIELD & 3 & send\_support \\
S3\_EMAIL\_SUBJECT & JSON\_FIELD & 3 & email\_subject \\
S3\_EMAIL\_BODY & JSON\_FIELD & 3 & email\_body \\
S4\_UPDATE & TOOL\_CALL & 3 & update\_result \\
S4\_EXEC\_RAW & TOOL\_CALL & 3 & exec\_raw \\
S4\_EXEC & TOOL\_CALL & 3 & exec\_result \\
S4\_BRANCH & RULE\_EVAL & 3 & support\_needed \\
S4\_SUPPORT\_RAW & TOOL\_CALL & 3 & support\_raw \\
S4\_SUPPORT & TOOL\_CALL & 3 & support\_result \\
S4\_VERIFY & TOOL\_CALL & 3 & score\_json \\
S5\_CHECK & RULE\_EVAL & 3 & passed \\
S6\_REPAIR\_1 & TOOL\_CALL & 3 & repair\_exec\_1 \\
S6\_REPAIR\_SUPPORT\_1 & TOOL\_CALL & 3 & repair\_support\_1 \\
S6\_RESCORE\_1 & TOOL\_CALL & 3 & score\_json\_1 \\
S6\_CHECK\_1 & RULE\_EVAL & 3 & passed\_1 \\
S6\_REPAIR\_2 & TOOL\_CALL & 3 & repair\_exec\_2 \\
S6\_REPAIR\_SUPPORT\_2 & TOOL\_CALL & 3 & repair\_support\_2 \\
S6\_RESCORE\_2 & TOOL\_CALL & 3 & score\_json\_2 \\
S6\_CHECK\_2 & RULE\_EVAL & 3 & passed\_2 \\
S7\_FINAL & FINAL\_ACTION & 3 & final\_result \\
S7\_AUDIT & AUDIT & 3 & auditJson \\
\midrule
\multicolumn{4}{l}{\textbf{FLOW\_AGENT\_OPS}: 37 节点 / 41 路由，operations（全确定性）} \\
\midrule
O\_START & AUDIT & 3 & ops\_audit\_start \\
O\_READ\_UNITS & TOOL\_CALL & 3 & ops\_ev\_units \\
O\_READ\_ENG & TOOL\_CALL & 3 & ops\_ev\_eng \\
O\_READ\_FUEL & TOOL\_CALL & 3 & ops\_ev\_fuel \\
O\_JOIN & AUDIT & 3 & ops\_audit\_join \\
O\_ROWS\_UNITS & ROWS\_JSON & 2 & ops\_units\_rows \\
O\_ROWS\_ENG & ROWS\_JSON & 2 & ops\_eng\_rows \\
O\_ROWS\_FUEL & ROWS\_JSON & 2 & ops\_fuel\_rows \\
O\_WHERE\_FUEL & WHERE\_JSON & 4 & ops\_rationed\_rows \\
O\_RATIONED\_BLDG & JSON\_FIELD & 3 & ops\_rationed\_building \\
O\_WHERE\_DUE & WHERE\_JSON & 4 & ops\_due\_rows \\
O\_DUE\_GATE & RULE\_EVAL & 3 & ops\_due\_ok \\
O\_U0\_ID & JSON\_FIELD & 3 & ops\_unit\_id\_0 \\
O\_U0\_BLDG & JSON\_FIELD & 3 & ops\_unit\_bldg\_0 \\
O\_EVENT0 & TOOL\_CALL & 3 & ops\_event\_0 \\
O\_ENG0 & WHERE\_JSON & 4 & ops\_eng\_0\_rows \\
O\_EMAIL0 & JSON\_FIELD & 3 & ops\_eng\_email\_0 \\
O\_MAILRAW0 & TOOL\_CALL & 3 & ops\_mail\_raw\_0 \\
O\_SEND0 & TOOL\_CALL & 3 & ops\_send\_0 \\
O\_U1\_ID & JSON\_FIELD & 3 & ops\_unit\_id\_1 \\
O\_U1\_BLDG & JSON\_FIELD & 3 & ops\_unit\_bldg\_1 \\
O\_EVENT1 & TOOL\_CALL & 3 & ops\_event\_1 \\
O\_ENG1 & WHERE\_JSON & 4 & ops\_eng\_1\_rows \\
O\_EMAIL1 & JSON\_FIELD & 3 & ops\_eng\_email\_1 \\
O\_MAILRAW1 & TOOL\_CALL & 3 & ops\_mail\_raw\_1 \\
O\_SEND1 & TOOL\_CALL & 3 & ops\_send\_1 \\
O\_U2\_ID & JSON\_FIELD & 3 & ops\_unit\_id\_2 \\
O\_U2\_BLDG & JSON\_FIELD & 3 & ops\_unit\_bldg\_2 \\
O\_EVENT2 & TOOL\_CALL & 3 & ops\_event\_2 \\
O\_ENG2 & WHERE\_JSON & 4 & ops\_eng\_2\_rows \\
O\_EMAIL2 & JSON\_FIELD & 3 & ops\_eng\_email\_2 \\
O\_MAILRAW2 & TOOL\_CALL & 3 & ops\_mail\_raw\_2 \\
O\_SEND2 & TOOL\_CALL & 3 & ops\_send\_2 \\
O\_VERIFY & TOOL\_CALL & 3 & ops\_score \\
O\_CHECK & RULE\_EVAL & 3 & ops\_passed \\
O\_FINAL & FINAL\_ACTION & 3 & ops\_final \\
O\_AUDIT & AUDIT & 3 & auditJson \\
\midrule
\multicolumn{4}{l}{\textbf{FLOW\_AGENT\_SUPPORT}: 60 节点 / 70 路由，support（1 次 LLM）} \\
\midrule
S\_START & AUDIT & 3 & support\_audit\_start \\
S\_READ\_BLOCK & TOOL\_CALL & 3 & support\_ev\_block \\
S\_READ\_SLA & TOOL\_CALL & 3 & support\_ev\_sla \\
S\_READ\_CONFIG & TOOL\_CALL & 3 & support\_ev\_config \\
S\_READ\_TICKETS & TOOL\_CALL & 3 & support\_ev\_tickets \\
S\_READ\_USERS & TOOL\_CALL & 3 & support\_ev\_users \\
S\_READ\_ORGS & TOOL\_CALL & 3 & support\_ev\_orgs \\
S\_READ\_CONTACTS & TOOL\_CALL & 3 & support\_ev\_contacts \\
S\_READ\_ACCOUNTS & TOOL\_CALL & 3 & support\_ev\_accounts \\
S\_READ\_CASES & TOOL\_CALL & 3 & support\_ev\_cases \\
S\_JOIN & AUDIT & 3 & support\_audit\_join \\
S\_EVIDENCE & JSON\_FIELD & 3 & support\_evidence \\
S\_FACTS & LLM\_CALL & 2 & support\_facts\_json \\
S\_CREATE\_ROWS & JSON\_FIELD & 3 & support\_create\_rows \\
S\_SHAPE\_GATE & RULE\_EVAL & 3 & support\_shape\_ok \\
S\_C0\_SUBJ & JSON\_FIELD & 3 & support\_c0\_subject \\
S\_C0\_ACC & JSON\_FIELD & 3 & support\_c0\_account\_id \\
S\_C0\_PRI & JSON\_FIELD & 3 & support\_c0\_priority \\
S\_C0\_CONTACT & JSON\_FIELD & 3 & support\_c0\_contact\_id \\
S\_C0\_NAME & JSON\_FIELD & 3 & support\_c0\_account\_name \\
S\_C0\_TID & JSON\_FIELD & 3 & support\_c0\_ticket\_id \\
S\_C0\_CREATE & TOOL\_CALL & 3 & support\_case\_0 \\
S\_C0\_COMMENT & TOOL\_CALL & 3 & support\_comment\_0 \\
S\_C1\_SUBJ & JSON\_FIELD & 3 & support\_c1\_subject \\
S\_C1\_ACC & JSON\_FIELD & 3 & support\_c1\_account\_id \\
S\_C1\_PRI & JSON\_FIELD & 3 & support\_c1\_priority \\
S\_C1\_CONTACT & JSON\_FIELD & 3 & support\_c1\_contact\_id \\
S\_C1\_NAME & JSON\_FIELD & 3 & support\_c1\_account\_name \\
S\_C1\_TID & JSON\_FIELD & 3 & support\_c1\_ticket\_id \\
S\_C1\_CREATE & TOOL\_CALL & 3 & support\_case\_1 \\
S\_C1\_COMMENT & TOOL\_CALL & 3 & support\_comment\_1 \\
S\_C2\_SUBJ & JSON\_FIELD & 3 & support\_c2\_subject \\
S\_C2\_ACC & JSON\_FIELD & 3 & support\_c2\_account\_id \\
S\_C2\_PRI & JSON\_FIELD & 3 & support\_c2\_priority \\
S\_C2\_CONTACT & JSON\_FIELD & 3 & support\_c2\_contact\_id \\
S\_C2\_NAME & JSON\_FIELD & 3 & support\_c2\_account\_name \\
S\_C2\_TID & JSON\_FIELD & 3 & support\_c2\_ticket\_id \\
S\_C2\_CREATE & TOOL\_CALL & 3 & support\_case\_2 \\
S\_C2\_COMMENT & TOOL\_CALL & 3 & support\_comment\_2 \\
S\_C3\_SUBJ & JSON\_FIELD & 3 & support\_c3\_subject \\
S\_C3\_ACC & JSON\_FIELD & 3 & support\_c3\_account\_id \\
S\_C3\_PRI & JSON\_FIELD & 3 & support\_c3\_priority \\
S\_C3\_CONTACT & JSON\_FIELD & 3 & support\_c3\_contact\_id \\
S\_C3\_NAME & JSON\_FIELD & 3 & support\_c3\_account\_name \\
S\_C3\_TID & JSON\_FIELD & 3 & support\_c3\_ticket\_id \\
S\_C3\_CREATE & TOOL\_CALL & 3 & support\_case\_3 \\
S\_C3\_COMMENT & TOOL\_CALL & 3 & support\_comment\_3 \\
S\_C4\_SUBJ & JSON\_FIELD & 3 & support\_c4\_subject \\
S\_C4\_ACC & JSON\_FIELD & 3 & support\_c4\_account\_id \\
S\_C4\_PRI & JSON\_FIELD & 3 & support\_c4\_priority \\
S\_C4\_CONTACT & JSON\_FIELD & 3 & support\_c4\_contact\_id \\
S\_C4\_NAME & JSON\_FIELD & 3 & support\_c4\_account\_name \\
S\_C4\_TID & JSON\_FIELD & 3 & support\_c4\_ticket\_id \\
S\_C4\_CREATE & TOOL\_CALL & 3 & support\_case\_4 \\
S\_C4\_COMMENT & TOOL\_CALL & 3 & support\_comment\_4 \\
S\_SLACK & TOOL\_CALL & 3 & support\_slack \\
S\_VERIFY & TOOL\_CALL & 3 & support\_score \\
S\_CHECK & RULE\_EVAL & 3 & support\_passed \\
S\_FINAL & FINAL\_ACTION & 3 & support\_final \\
S\_AUDIT & AUDIT & 3 & auditJson \\
\midrule
\multicolumn{4}{l}{\textbf{FLOW\_AGENT\_ULTRA\_NOREPAIR}: 28 节点 / 34 路由，消融：去掉 repair} \\
\midrule
S0\_CONTRACT & LLM\_CALL & 2 & contract\_json \\
S1\_FORK & AUDIT & 3 & audit\_s1\_start \\
S1\_GMAIL & TOOL\_CALL & 3 & ev\_gmail \\
S1\_HIER & TOOL\_CALL & 3 & ev\_hier \\
S1\_FX & TOOL\_CALL & 3 & ev\_fx \\
S1\_OPP & TOOL\_CALL & 3 & ev\_opp \\
S1\_CASES & TOOL\_CALL & 3 & ev\_cases \\
S1\_JOIN & AUDIT & 3 & audit\_s1\_join \\
S1\_CONTEXT & CONTEXT\_JSON & 1 & evidence\_ctx \\
S2\_FACTS & LLM\_CALL & 2 & facts\_json \\
S2\_CHECK & RULE\_EVAL & 3 & facts\_valid \\
S3\_PLAN & LLM\_CALL & 2 & plan\_json \\
S3\_UPDATE\_STAGE & JSON\_FIELD & 3 & update\_stage \\
S3\_EXEC\_TO & JSON\_FIELD & 3 & exec\_to \\
S3\_SUPPORT\_TO & JSON\_FIELD & 3 & support\_to \\
S3\_SEND\_FLAG & JSON\_FIELD & 3 & send\_support \\
S3\_EMAIL\_SUBJECT & JSON\_FIELD & 3 & email\_subject \\
S3\_EMAIL\_BODY & JSON\_FIELD & 3 & email\_body \\
S4\_UPDATE & TOOL\_CALL & 3 & update\_result \\
S4\_EXEC\_RAW & TOOL\_CALL & 3 & exec\_raw \\
S4\_EXEC & TOOL\_CALL & 3 & exec\_result \\
S4\_BRANCH & RULE\_EVAL & 3 & support\_needed \\
S4\_SUPPORT\_RAW & TOOL\_CALL & 3 & support\_raw \\
S4\_SUPPORT & TOOL\_CALL & 3 & support\_result \\
S4\_VERIFY & TOOL\_CALL & 3 & score\_json \\
S5\_CHECK & RULE\_EVAL & 3 & passed \\
S7\_FINAL & FINAL\_ACTION & 3 & final\_result \\
S7\_AUDIT & AUDIT & 3 & auditJson \\
\midrule
\multicolumn{4}{l}{\textbf{FLOW\_AGENT\_OPS\_SEQ}: 37 节点 / 39 路由，消融：串行读取} \\
\midrule
O\_START & AUDIT & 3 & ops\_audit\_start \\
O\_READ\_UNITS & TOOL\_CALL & 3 & ops\_ev\_units \\
O\_READ\_ENG & TOOL\_CALL & 3 & ops\_ev\_eng \\
O\_READ\_FUEL & TOOL\_CALL & 3 & ops\_ev\_fuel \\
O\_JOIN & AUDIT & 3 & ops\_audit\_join \\
O\_ROWS\_UNITS & ROWS\_JSON & 2 & ops\_units\_rows \\
O\_ROWS\_ENG & ROWS\_JSON & 2 & ops\_eng\_rows \\
O\_ROWS\_FUEL & ROWS\_JSON & 2 & ops\_fuel\_rows \\
O\_WHERE\_FUEL & WHERE\_JSON & 4 & ops\_rationed\_rows \\
O\_RATIONED\_BLDG & JSON\_FIELD & 3 & ops\_rationed\_building \\
O\_WHERE\_DUE & WHERE\_JSON & 4 & ops\_due\_rows \\
O\_DUE\_GATE & RULE\_EVAL & 3 & ops\_due\_ok \\
O\_U0\_ID & JSON\_FIELD & 3 & ops\_unit\_id\_0 \\
O\_U0\_BLDG & JSON\_FIELD & 3 & ops\_unit\_bldg\_0 \\
O\_EVENT0 & TOOL\_CALL & 3 & ops\_event\_0 \\
O\_ENG0 & WHERE\_JSON & 4 & ops\_eng\_0\_rows \\
O\_EMAIL0 & JSON\_FIELD & 3 & ops\_eng\_email\_0 \\
O\_MAILRAW0 & TOOL\_CALL & 3 & ops\_mail\_raw\_0 \\
O\_SEND0 & TOOL\_CALL & 3 & ops\_send\_0 \\
O\_U1\_ID & JSON\_FIELD & 3 & ops\_unit\_id\_1 \\
O\_U1\_BLDG & JSON\_FIELD & 3 & ops\_unit\_bldg\_1 \\
O\_EVENT1 & TOOL\_CALL & 3 & ops\_event\_1 \\
O\_ENG1 & WHERE\_JSON & 4 & ops\_eng\_1\_rows \\
O\_EMAIL1 & JSON\_FIELD & 3 & ops\_eng\_email\_1 \\
O\_MAILRAW1 & TOOL\_CALL & 3 & ops\_mail\_raw\_1 \\
O\_SEND1 & TOOL\_CALL & 3 & ops\_send\_1 \\
O\_U2\_ID & JSON\_FIELD & 3 & ops\_unit\_id\_2 \\
O\_U2\_BLDG & JSON\_FIELD & 3 & ops\_unit\_bldg\_2 \\
O\_EVENT2 & TOOL\_CALL & 3 & ops\_event\_2 \\
O\_ENG2 & WHERE\_JSON & 4 & ops\_eng\_2\_rows \\
O\_EMAIL2 & JSON\_FIELD & 3 & ops\_eng\_email\_2 \\
O\_MAILRAW2 & TOOL\_CALL & 3 & ops\_mail\_raw\_2 \\
O\_SEND2 & TOOL\_CALL & 3 & ops\_send\_2 \\
O\_VERIFY & TOOL\_CALL & 3 & ops\_score \\
O\_CHECK & RULE\_EVAL & 3 & ops\_passed \\
O\_FINAL & FINAL\_ACTION & 3 & ops\_final \\
O\_AUDIT & AUDIT & 3 & auditJson \\
\end{longtable}
}


\section{逐断言结果（direct vs structured）}

\noindent{\footnotesize P = 通过，F = 失败；只列非 excluded 断言。\par}


\subsection*{sales.multi\_hop\_lookup}

{\footnotesize
\begin{longtable}{p{1.5cm}p{1.5cm}p{10.8cm}}
\toprule
direct & structured & 断言（类型 + 关键参数） \\
\midrule
\endhead
\bottomrule
\endfoot
\TiiMark{P} & \TiiMark{P} & salesforce\_field\_equals\newline collection=opportunities \\
\TiiMark{F} & \TiiMark{P} & gmail\_message\_sent\_to\_with\_body\_contains\newline subject=Deal Closed Notification; to=support-escalation@example.com; body\_contains=Meridian Corp - Platform Deal/\$156,000/Enterprise \\
\TiiMark{F} & \TiiMark{P} & gmail\_message\_sent\_to\_with\_body\_contains\newline subject=Deal Closed Notification; to=executive-team@example.com; body\_contains=Meridian Corp - Platform Deal/\$156,000/Enterprise \\
\end{longtable}
}


\subsection*{operations.generator\_load\_testing\_schedule}

{\footnotesize
\begin{longtable}{p{1.5cm}p{1.5cm}p{10.8cm}}
\toprule
direct & structured & 断言（类型 + 关键参数） \\
\midrule
\endhead
\bottomrule
\endfoot
\TiiMark{P} & \TiiMark{P} & google\_calendar\_event\_exists\newline summary\_contains=Generator Load Test - GEN-01 \\
\TiiMark{P} & \TiiMark{P} & google\_calendar\_event\_exists\newline summary\_contains=Generator Load Test - GEN-02 \\
\TiiMark{P} & \TiiMark{P} & google\_calendar\_event\_exists\newline summary\_contains=Generator Load Test - GEN-07 \\
\TiiMark{P} & \TiiMark{P} & gmail\_message\_sent\_to\_with\_body\_contains\newline to=c.ruiz@company.example.com; body\_contains=GEN-01 \\
\TiiMark{P} & \TiiMark{P} & gmail\_message\_sent\_to\_with\_body\_contains\newline to=d.west@company.example.com; body\_contains=GEN-02 \\
\TiiMark{P} & \TiiMark{P} & gmail\_message\_sent\_to\_with\_body\_contains\newline to=e.popov@company.example.com; body\_contains=GEN-07 \\
\TiiMark{F} & \TiiMark{P} & gmail\_message\_sent\newline to=c.ruiz@company.example.com; body\_contains=3 \\
\end{longtable}
}


\subsection*{support.zendesk\_sf\_case\_sync}

{\footnotesize
\begin{longtable}{p{1.5cm}p{1.5cm}p{10.8cm}}
\toprule
direct & structured & 断言（类型 + 关键参数） \\
\midrule
\endhead
\bottomrule
\endfoot
\TiiMark{P} & \TiiMark{P} & salesforce\_case\_exists\newline subject=Cannot access reporting dashboard; origin=Web; account\_id=sf\_a1 \\
\TiiMark{F} & \TiiMark{P} & salesforce\_case\_exists\newline subject=Mobile app slow on enterprise plan; priority=High; origin=Web; account\_id=sf\_a3 \\
\TiiMark{P} & \TiiMark{P} & salesforce\_case\_exists\newline subject=Data import issues; priority=Critical; origin=Web; account\_id=sf\_a1 \\
\TiiMark{P} & \TiiMark{P} & salesforce\_case\_exists\newline subject=Webhook delivery delays; priority=Low; origin=Web; account\_id=sf\_a4 \\
\TiiMark{P} & \TiiMark{P} & zendesk\_ticket\_has\_comment\newline ticket\_id=tkt\_101; body\_contains=Salesforce \\
\TiiMark{P} & \TiiMark{P} & zendesk\_ticket\_has\_comment\newline ticket\_id=tkt\_102; body\_contains=Salesforce \\
\TiiMark{P} & \TiiMark{P} & zendesk\_ticket\_has\_comment\newline ticket\_id=tkt\_107; body\_contains=Salesforce \\
\TiiMark{P} & \TiiMark{P} & zendesk\_ticket\_has\_comment\newline ticket\_id=tkt\_110; body\_contains=Salesforce \\
\TiiMark{P} & \TiiMark{P} & zendesk\_ticket\_has\_comment\newline ticket\_id=tkt\_111; body\_contains=Salesforce \\
\TiiMark{F} & \TiiMark{P} & salesforce\_collection\_count\_equals\newline count=6; collection=cases \\
\TiiMark{F} & \TiiMark{P} & slack\_message\_exists\newline channel\_name=support-sync; text\_contains=5 \\
\TiiMark{F} & \TiiMark{P} & slack\_message\_exists\newline channel\_name=support-sync; text\_contains=6 \\
\TiiMark{P} & \TiiMark{P} & slack\_message\_exists\newline channel\_name=support-sync; text\_contains=Critical \\
\TiiMark{P} & \TiiMark{P} & slack\_message\_in\_channel\newline channel\_name=support-sync; text\_contains=SYNC-OPS-20260210 \\
\TiiMark{P} & \TiiMark{P} & zendesk\_ticket\_has\_comment\newline ticket\_id=tkt\_102; body\_contains=NovaTech \\
\TiiMark{P} & \TiiMark{P} & slack\_message\_in\_channel\newline channel\_name=support-sync; text\_contains=Acme \\
\end{longtable}
}


\end{document}
